\begin{abstract}
本文针对多源融合机器人定位及任务优化问题，建立了从低阶到高阶、从单一偏差到多源耦合的完整建模体系；建立了“时间对齐—偏差识别—异步融合—任务调度”的统一框架。

针对问题一，分别采用三次样条构造两种定位方式的连续轨迹，基于速度特征的归一化互相关对时间偏差进行全局粗估计，再以位置均方误差为目标，采用 Brent 法完成一维精细搜索。得到方式2时间戳的加性修正量为 $-198.4317$ s，即其时间轴前移 $198.4317$ s。合并两路无噪声观测后采用三次样条重建，生成 $221.0\sim1070.4$ s 的 8495 个 10 Hz 轨迹点；700 个同刻观测坐标完全一致，对齐 RMSE 为 $8.51\times10^{-11}$ m。

对于问题二，在固定公共网格上以逐轴 Huber 剖面损失联合估计时间与空间相对偏差，再用异步常加速度 Kalman 滤波和 RTS 固定区间平滑融合原始事件。估得 $t_{2,\mathrm{ref}}=t_2-50.171732$ s，相对系统偏差为 $(3.466998,-1.833200)$ m。过程噪声由无门控的预检验 NIS 与白化创新相关性选择，最终输出 $102.0\sim951.5$ s 的 8496 个 10 Hz 状态点。

对于问题三，先对实际测量数据作稳健时间配准，再以 HAC-Wald 检验、工程效应阈值及移动区块自助法联合判断系统偏差。估得 $\widehat{\Delta t}=-367.88$ s；Wald 检验 $p=0.3511$，工程效应指数为 $0.0681$，两轴自助置信区间均包含0，故未发现需要启用偏差状态的证据。随后完成选择性异步融合，得到 $445.0076\sim814.0076$ s 的 3691 个 10 Hz 轨迹点，并报告位置不确定性。

对于问题四，将距离、速度、加速度和完整准备时间转化为连续可行窗口，经 0.01 s 加密复核后构造 1255 个候选，以三阶段词典序混合整数规划依次最大化任务数、最小约束裕度和综合质量。最终完成模板上限的9项任务，其中射击4项、拍照5项，最小归一化裕度为0.4241，射击期望命中数为3.40；三档各500次扰动下整套方案均保持可行。结果已写入独立的 \texttt{result.xlsx} 副本，原始模板保持不变。

\keywords{多源定位\quad 时间配准\quad Huber估计\quad Kalman滤波\quad 混合整数规划}
\end{abstract}
